\documentclass[11pt]{article}
\usepackage[margin=1in]{geometry}
\usepackage{amsmath,amssymb,amsthm,mathtools}
\usepackage{graphicx}
\usepackage{microtype}
\usepackage{xcolor}
\usepackage[hidelinks]{hyperref}
\usepackage{enumitem}

\newtheorem{theorem}{Theorem}
\newtheorem{proposition}[theorem]{Proposition}
\newtheorem{lemma}[theorem]{Lemma}
\newtheorem{corollary}[theorem]{Corollary}
\theoremstyle{remark}
\newtheorem{remark}[theorem]{Remark}

\newcommand{\M}{\mathcal M}
\newcommand{\N}{\mathcal N}
\newcommand{\tr}{\operatorname{tr}}
\newcommand{\id}{\operatorname{id}}
\newcommand{\Span}{\operatorname{span}}
\newcommand{\C}{\mathbb C}
\newcommand{\Nn}{\mathbb N}
\newcommand{\one}{\mathbf 1}
\newcommand{\Spol}{\mathfrak S}

\title{Finite matrix levels in noncommutative $L_p$:\\
trace-identity lower bounds at even exponents}
\author{Discovery packet for arXiv:0707.0427}
\date{17 August 2026}

\begin{document}
\maketitle

\begin{center}
\fbox{\parbox{0.92\linewidth}{\textbf{Status: substantial partial result,
likely valid.}  For $p=2m$ this packet constructs finite-dimensional maps
which agree exactly through a prescribed matrix level and fail at the next
one.  The non-even finiteness questions remain open.}}
\end{center}

\section{Source question and result}

The source is Mikael de la Salle, \emph{Complete isometries between subspaces
of noncommutative $L_p$-spaces}, arXiv:0707.0427, published in the
\emph{Journal of Operator Theory} \cite{dls}.  Section~3.3 (printed page~24)
asks for two universal matrix levels.  In the finite tracial setting,
$n_{p,1}$ is the least $n$ such that
\begin{equation}\label{eq:unit-data}
 \|\one_n\otimes\one+X\|_p
 =\|\one_n\otimes\one+(\id_n\otimes u)X\|_p
 \quad(X\in M_n(E))
\end{equation}
forces the same equality at every matrix size.  In the semifinite tracial
setting, $n_p$ is the least $n$ such that every $n$-isometry between
subspaces is completely isometric.  The source proves
\begin{equation}\label{eq:source-upper}
 n_{2m,1}\le m,\qquad n_{2m}\le m,
\end{equation}
and records $n_{p,1},n_p>1$ for $p\ne2$, using transposition.  It leaves open
whether either invariant is finite when $p$ is not even.

\begin{figure}[ht]
\centering
\includegraphics[width=0.94\linewidth]{figures/open_problem_crop.png}
\caption{The exact questions and the source upper bounds, arXiv:0707.0427,
Section~3.3, printed page~24.}
\end{figure}

The main result supplies quantitative lower bounds in the even case.

\begin{theorem}[Matrix-level lower bounds]\label{thm:main}
Let $m\ge2$ be an integer.
\begin{enumerate}[label=\textup{(\roman*)}]
\item For every $1\le r\le\lfloor m/2\rfloor$ there are finite-dimensional
tracial von Neumann algebras $(\M,\tau)$ and $(\N,\widetilde\tau)$,
finite-dimensional subspaces $E\subset\M$, $F\subset\N$, and a linear map
$u:E\to F$ satisfying \eqref{eq:unit-data} at size $r$ but not at size
$r+1$.
\item For every $1\le r\le\lfloor(m-1)/2\rfloor$ there are such finite
tracial algebras, subspaces, and a linear map $u:E\to F$ which is
$r$-isometric but not $(r+1)$-isometric.
\end{enumerate}
Consequently
\begin{equation}\label{eq:bounds}
 \boxed{\ \lfloor m/2\rfloor+1\le n_{2m,1}\le m\ },
 \qquad
 \boxed{\ \max\{2,\lceil m/2\rceil\}\le n_{2m}\le m\ }.
\end{equation}
All counterexamples in the lower-bound statements can be taken between
subspaces of finite direct sums of full matrix algebras, with faithful
normalized traces and bounded, linearly independent spanning tuples.
\end{theorem}

In particular $n_{4,1}=n_4=2$.  The new lower bound for $n_{2m,1}$ improves
the transposition bound from $p=8$ onward; the new lower bound for $n_{2m}$
improves it from $p=10$ onward.

\section{Proof intuition}

At an even exponent, the $2m$th power of an $L_{2m}$ norm is a finite cyclic
$*$-polynomial.  Fixed size $r$ sees that polynomial only through traces of
$r\times r$ coefficient matrices.  The Amitsur--Levitzki identity makes the
alternating polynomial in $2r$ matrices vanish identically at that size, but
not at size $r+1$.  Appending one distinguished marker prevents the
alternating sum from disappearing merely by cyclicity of the trace.  This
exhibits a direction in the level-$(r+1)$ norm data which is invisible at
level $r$.

The remaining issue is positivity: a formal difference of moments need not
come from two tracial probability spaces.  Because only finitely many moment
degrees occur at $p=2m$, finite-dimensional separation resolves this.  Matrix
trace evaluations span the dual of the relevant cyclic-polynomial space.
Write a separating functional as a signed combination of evaluations, split
it into positive and negative parts, and realize those convex combinations
as traces on finite direct sums of matrix algebras.  A common tagging summand
makes the two spanning tuples linearly independent.  This converts the trace
identity into an honest linear map of $L_{2m}$ subspaces.

For non-even $p$, infinitely many moment degrees occur.  The same finite
separation no longer matches all coefficients, which is precisely why the
original non-even problem is not resolved here.

\section{Preliminaries}

For matrices $B_1,\ldots,B_q$, write
\[
 \Spol_q(B_1,\ldots,B_q)
 =\sum_{\sigma\in S_q}\operatorname{sgn}(\sigma)
   B_{\sigma(1)}\cdots B_{\sigma(q)}.
\]
The Amitsur--Levitzki theorem says that $\Spol_{2r}$ vanishes identically on
$M_r(\C)$ \cite{AL}.  The following explicit evaluation also records the
sharpness needed below.

\begin{lemma}[Chain witness]\label{lem:chain}
In $M_{r+1}$ put
\[
 B_1=e_{11},\ B_2=e_{12},\ B_3=e_{22},\ B_4=e_{23},\ldots,
 B_{2r-1}=e_{rr},\ B_{2r}=e_{r,r+1}.
\]
Then
\[
 \Spol_{2r}(B_1,\ldots,B_{2r})=e_{1,r+1}.
\]
In particular, with $B_0=e_{r+1,1}$,
$\tr_{r+1}(\Spol_{2r}(B_1,\ldots,B_{2r})B_0)=1/(r+1)$.
\end{lemma}

\begin{proof}
The displayed ordering multiplies to $e_{1,r+1}$.  Any nonzero product of
all the listed matrix units must follow the directed chain and must place
each vertex loop immediately before the outgoing edge at that vertex.  Thus
the displayed ordering is the unique nonzero summand of the standard
polynomial, and it has positive sign.
\end{proof}

We also use an elementary finite-dimensional realization principle.  A
cyclic $*$-polynomial means a noncommutative $*$-polynomial modulo cyclic
rotation of words; evaluation is always followed by a trace.

\begin{lemma}[Finite tracial separation]\label{lem:separation}
Let $V$ be a finite-dimensional real space of self-adjoint cyclic
$*$-polynomials in finitely many variables.  The real linear span of the
functionals
\[
 q\longmapsto \tr_d(q(Z_1,\ldots,Z_k)),
 \qquad d\in\Nn,\quad Z_i\in M_d,
\]
is $V^*$.

Suppose $U\subset V$ contains the constant polynomial and $q_0\notin U$.
Then there are finite-dimensional tracial von Neumann algebras and tuples
$x,y$ such that their tracial evaluations agree on $U$ and disagree on
$q_0$.
\end{lemma}

\begin{proof}
Finite matrix traces separate nonzero cyclic polynomials.  One elementary
way to see this is to choose a cyclic word with nonzero coefficient, place
independent scalar-weighted matrix units around a directed cycle spelling
that word, and read its scalar monomial in the resulting trace.  Polarizing
separates repeated letters and taking the adjoint handles the self-adjoint
real form.  Hence the displayed evaluation functionals have zero common
annihilator in $V$, so they span $V^*$.

Choose $\Lambda\in V^*$ with $\Lambda|_U=0$ and
$\Lambda(q_0)\ne0$, and write
$\Lambda=\sum_j c_j\operatorname{ev}_{Z^{(j)}}$.  Since $1\in U$,
$\sum_jc_j=0$.  The positive and negative coefficient sums therefore have
the same value $T>0$.  Dividing by $T$, each side is a convex combination of
normalized matrix traces.  Such a combination is the faithful normalized
trace on the direct sum of the corresponding matrix algebras, with the
coefficients as positive central weights.  The resulting direct-sum tuples
$x$ and $y$ have moment-functional difference $\Lambda/T$.
\end{proof}

If linear independence of the tuples is needed, replace both by
$x_i\oplus t_i$ and $y_i\oplus t_i$, use the same positive weight on the
common summand, and choose linearly independent matrix units $t_i$.  This
multiplies their moment-functional difference by a nonzero scalar and makes
the prescribed basis map well defined.

\section{An optimal cyclic-isolation lemma}

The source constructs $N$ matrices of size $\lceil N/2\rceil$ whose trace
selects exactly the cyclic order among all permutations.  The next
proposition shows that its matrix size cannot be improved.

\begin{proposition}[Exact isolation size]\label{prop:isolation}
Let $c(N)$ be the least $d$ for which there exist $a_1,\ldots,a_N\in M_d$
such that
\[
 \tr_d(a_{\sigma(1)}\cdots a_{\sigma(N)})=
 \begin{cases}
  1,&\sigma\text{ is a cyclic rotation of }(1,\ldots,N),\\
  0,&\text{otherwise}.
 \end{cases}
\]
Then $c(N)=\lceil N/2\rceil$.
\end{proposition}

\begin{proof}
The upper bound is the matrix-unit construction in Remark~1.5 of the source
\cite{dls}.  For the reverse bound, suppose $d<\lceil N/2\rceil$.  Then
$N-2d\ge1$.  Amitsur--Levitzki gives
\[
 0=\tr_d\!\left(
 \Spol_{2d}(a_1,\ldots,a_{2d})a_{2d+1}\cdots a_N\right).
\]
In the alternating expansion, the full ordering
$\sigma(1),\ldots,\sigma(2d),2d+1,\ldots,N$ is a cyclic rotation of
$1,\ldots,N$ only for the identity permutation: its last entry is $N$, so a
cyclic rotation must start at $1$.  The assumed values therefore make the
same alternating sum equal to one, a contradiction.
\end{proof}

Combining this with Lemma~2.2 and Lemma~3.2 of the source gives a useful
positive statement: level $r$ unit-perturbation equality recovers a cyclic
word of degree $N\le2r$ whenever the source's binomial coefficient is
nonzero.  In particular, for non-even $p$ it recovers every $*$-moment of
degree at most $2r$.  Proposition~\ref{prop:isolation} says the direct
cyclic-selector method has an exact degree horizon $2r$.

\section{Proof of the unit-perturbation lower bound}

Fix $m\ge2$ and $r$ with $2r\le m$.  Work with formal variables
$X_0,X_1,\ldots,X_{2r}$.  For $s\ge1$ and coefficient matrices
$A=(A_0,\ldots,A_{2r})\in M_s^{2r+1}$, let $P_{s,A}$ be the cyclic
$*$-polynomial characterized by
\begin{equation}\label{eq:P}
 \tau(P_{s,A}(x))
 =\left\|\one_s\otimes\one+\sum_{j=0}^{2r}A_j\otimes x_j
  \right\|_{2m}^{2m}.
\end{equation}
It has degree at most $2m$.  Let $U_s$ be the real span of the self-adjoint
real and imaginary parts of all $P_{s,A}$.

Give exactly one of $X_1,\ldots,X_{2r}$ a star and leave all other variables,
including $X_0$, unstarred.  Define a coefficient functional $L$ on cyclic
polynomials by alternating the coefficients of the cyclic words
\[
 X_{\sigma(1)}^{\varepsilon_{\sigma(1)}}\cdots
 X_{\sigma(2r)}^{\varepsilon_{\sigma(2r)}}X_0,
 \qquad \sigma\in S_{2r},
\]
with sign $\operatorname{sgn}(\sigma)$, and extend it to the self-adjoint
real form by taking the corresponding real or imaginary part.

Every displayed word has length $k=2r+1$ and exactly one cyclic
$*$-to-$1$ transition.  The coefficient formula in Lemma~2.2 of the source
therefore gives, up to a fixed nonzero cyclic-normalization factor,
\begin{equation}\label{eq:LPA}
 L(P_{s,A})=
 c_{m,r}\,\tr_s\!\left(
 \Spol_{2r}(B_1,\ldots,B_{2r})B_0\right),
\end{equation}
where $B_j=A_j^{\varepsilon_j}$, $B_0=A_0$, and
\begin{equation}\label{eq:c}
 c_{m,r}=(2r+1)\binom{m}{2r+1}+2r\binom{m}{2r}.
\end{equation}
This is nonzero for $m\ge2r$; at the boundary $m=2r$ it equals $2r$.

By Amitsur--Levitzki, \eqref{eq:LPA} is zero for $s=r$, so $L$ annihilates
$U_r$.  For $s=r+1$, choose the chain matrices from
Lemma~\ref{lem:chain}, take $B_0=e_{r+1,1}$, and choose the corresponding
$A_j$ (replace the uniquely starred coefficient by its adjoint).  Then
\eqref{eq:LPA} is nonzero.  Hence this particular $P_{r+1,A}$ does not lie
in $U_r$.

Apply Lemma~\ref{lem:separation} in the finite-dimensional cyclic-polynomial
space of degree at most $2m$, with $U=U_r$ and $q_0=P_{r+1,A}$.  We obtain
finite tracial tuples $x,y$ for which all level-$r$ quantities
\eqref{eq:P} agree, while this level-$(r+1)$ quantity differs.  Add a common
tagging summand to make each tuple linearly independent, set
$E=\Span\{x_j\}$, $F=\Span\{y_j\}$, and define $u(x_j)=y_j$.  Equality of
the $2m$th powers is equivalent to equality of the norms, so $u$ satisfies
the desired level-$r$ unit-perturbation equality and fails it at level
$r+1$.  This proves Theorem~\ref{thm:main}(i).

\section{Proof of the ordinary-isometry lower bound}

Now assume $2r+1\le m$.  Introduce pairs of variables
$(U_j,V_j)$ for $0\le j\le2r$ and, if
$h=m-(2r+1)>0$, filler pairs $(R_\ell,S_\ell)$ for $1\le\ell\le h$.
For a coefficient tuple $A$ define the homogeneous cyclic polynomial
$Q_{s,A}$ by
\begin{equation}\label{eq:Q}
 \tau(Q_{s,A}(z))=
 \left\|\sum_i A_i\otimes z_i\right\|_{2m}^{2m}
 = (\tr_s\otimes\tau)((Z^*Z)^m).
\end{equation}

Define $L_0$ by alternating the coefficients of the words
\[
 (U_{\sigma(1)}^*V_{\sigma(1)})\cdots
 (U_{\sigma(2r)}^*V_{\sigma(2r)})
 (U_0^*V_0)(R_1^*S_1)\cdots(R_h^*S_h).
\]
Expansion of $(Z^*Z)^m$ gives, up to one fixed positive cyclic factor,
\begin{equation}\label{eq:L0Q}
 L_0(Q_{s,A})=
 \tr_s\!\left(\Spol_{2r}(B_1,\ldots,B_{2r})B_0C_1\cdots C_h\right),
\end{equation}
where $B_j=A_{U_j}^*A_{V_j}$, $B_0=A_{U_0}^*A_{V_0}$, and
$C_\ell=A_{R_\ell}^*A_{S_\ell}$.

Equation~\eqref{eq:L0Q} vanishes identically at size $r$.  At size $r+1$,
take $A_{U_j}=\one$, choose $A_{V_j}$ to be the chain matrices, choose the
marker product to be $e_{r+1,1}$, and take every filler product to be the
identity.  Lemma~\ref{lem:chain} makes \eqref{eq:L0Q} nonzero.  Thus a
level-$(r+1)$ homogeneous norm polynomial lies outside the span of the
level-$r$ ones and the constant polynomial.

Apply Lemma~\ref{lem:separation} to that enlarged level-$r$ span.  The same
positive/negative direct-sum construction and common tagging summand produce
a well-defined linear map whose $M_r$ amplification is isometric and whose
$M_{r+1}$ amplification is not.  This proves Theorem~\ref{thm:main}(ii).

Taking all admissible $r$ gives
$n_{2m,1}\ge\lfloor m/2\rfloor+1$ and
$n_{2m}\ge\lceil m/2\rceil$.  Transposition supplies
$n_{2m}\ge2$, while \eqref{eq:source-upper} supplies both upper bounds in
\eqref{eq:bounds}.

\section{Verification, limitations, and novelty status}

The script \texttt{code/verify\_trace\_identity.py} performs exact integer
checks for $r=1,\ldots,4$: sampled instances of
$\Spol_{2r}=0$ in $M_r$, the explicit chain identity in $M_{r+1}$, the
marker trace, and $c_{2r,r}=2r$.  These checks audit the algebra but are not
used as proof.

The construction is existential after the trace-polynomial separation step;
it does not optimize the dimensions of the direct-sum algebras.  The bounds
do not determine either invariant except at $p=2$ and $p=4$.  Most
importantly, they do not settle finiteness for non-even $p$.  There the norm
expansion has infinitely many moment degrees, while Lemma~\ref{lem:separation}
only realizes a separator on a finite-dimensional truncation.  Producing two
positive tracial distributions agreeing on the entire fixed-level data space
is the remaining obstruction.

A bounded novelty search on 17 August 2026 checked the run indexes, the exact
arXiv id and notation $n_{p,1}$, the section title, and primary-source queries
combining fixed matrix level, $n$-isometry, complete isometry, subspaces, and
noncommutative $L_p$.  It found the source, the author's thesis, and general
whole-space isometry classifications, but no later answer or matching
trace-identity lower bound.  This is bounded evidence, not a priority claim.

\paragraph{Human-review recommendation.}
High-value specialist review is recommended.  The main gates are the
finite-tracial-separation lemma and the bookkeeping in
\eqref{eq:LPA} and \eqref{eq:L0Q}; the Amitsur--Levitzki witnesses themselves
are exact and computationally audited.

\begin{thebibliography}{9}
\bibitem{dls}
M.~de la Salle,
\emph{Complete isometries between subspaces of noncommutative $L_p$-spaces},
J. Operator Theory \textbf{64} (2010), 261--298;
arXiv:0707.0427.

\bibitem{AL}
A.~S. Amitsur and J.~Levitzki,
\emph{Minimal identities for algebras},
Proc. Amer. Math. Soc. \textbf{1} (1950), 449--463.
\end{thebibliography}

\end{document}

